%% ──────────────────────────────────────────────
%%  فصل ۱۵ – پارادوکس رهایی: وقتی نجات به انقیاد بدل می‌شود
%%  فایل: chapters/ch15-paradox-of-liberation.tex
%% ──────────────────────────────────────────────
\chapter{پارادوکس رهایی: وقتی نجات به انقیاد بدل می‌شود}
\label{ch:paradox-of-liberation}

\begin{tcolorbox}[mybox, title={درباره‌ی این فصل}]
این فصل — بر خلاف فصل‌های قبلی که تلاش می‌کردند توصیفی و تحلیلی باشند —
\textbf{نقدی صریح و چندوجهی} است. موضوع: چرا مشروعیت‌بخشی به مداخله‌ی خارجی
— حتی با بهترین نیت‌ها — باز کردن جعبه‌ی پاندوراست، و در عین حال چرا نمی‌توان
این بحث را به سادگی «همه‌چیز بد است» خلاصه کرد.
\end{tcolorbox}


%% ================================================================
\section{آمار تلخ: ۸۴٪ شکست}
\label{sec:ch15-failure-stats}
%% ================================================================

\subsection{داده‌ها چه می‌گویند؟}

\needspace{6\baselineskip}
پژوهش‌های جامع در مورد نتایج مداخلات خارجیِ «رهایی‌بخش» نتایج
تکان‌دهنده‌ای ارائه می‌دهند.  مؤسسه‌ی \lat{RAND}، برنامه‌ی داده‌های
درگیری اوپسالا (\lat{UCDP})، و گروه کاری بی‌ثباتی سیاسی
(\lat{Political Instability Task Force}) هر سه به الگوی مشابهی
رسیده‌اند:\footnote{Peksen, Dursun. ``When Do Imposed Sanctions Work?''
\textit{Defence and Peace Economics} 30, no.\,6 (2019): 635--647; Downes,
Alexander B. \textit{Catastrophic Success: Why Foreign-Imposed Regime
Change Goes Wrong}. Cornell UP, 2021.}

\begin{tcolorbox}[failbox, title={آمار شکست مداخلات رهایی‌بخش}]
\begin{table}[H]
\centering
\begin{tabular}{r|c}
\toprule
\textbf{شاخص} & \textbf{درصد تقریبی} \\
\midrule
مداخلاتی که به \textbf{بی‌ثباتی بلندمدت} منجر شدند & $\sim$۷۰٪ \\
مداخلاتی که \textbf{وضعیت حقوق بشر را بدتر} کردند & $\sim$۵۵٪ \\
مداخلاتی که به \textbf{جنگ داخلی جدید} منجر شدند & $\sim$۴۵٪ \\
مداخلاتی با \textbf{نتایج نسبتاً مثبت} (۵ سال پس از مداخله) & $\sim$۱۶٪ \\
مداخلاتی با \textbf{نتایج مثبت پایدار} (۲۰ سال پس از مداخله) & $<$۱۰٪ \\
\bottomrule
\end{tabular}
\end{table}
\end{tcolorbox}

این آمار به ما می‌گوید که \textbf{از هر شش تلاش رهایی‌بخش با کمک
خارجی، فقط یکی به نتیجه‌ی نسبتاً قابل‌قبول می‌رسد} — و حتی این
«موفقیت» اغلب موقتی است.  مقایسه‌ای ساده: اگر پزشکی عمل جراحی‌اش ۸۴٪
شکست می‌خورد، مجوزش لغو می‌شد.  اما در عرصه‌ی بین‌المللی، «جرّاحان»
مداخله‌گر بارها و بارها اجازه‌ی تکرار می‌یابند.


\subsection{چرا این نرخ شکست چنین بالاست؟}

\needspace{6\baselineskip}
چهار عامل ساختاری — نه اتفاقی — این نرخ بالا را توضیح می‌دهند:

\begin{tcolorbox}[enemybox, title={عامل ۱: عدم شناخت بافت محلی}]
مداخله‌گران معمولاً:
\begin{itemize}
\item زبان محلی را نمی‌دانند — در عراق، کمتر از ۱٪ کارکنان
\lat{CPA} عربی می‌دانستند\footnote{Chandrasekaran, Rajiv.
\textit{Imperial Life in the Emerald City}. Knopf, 2006.}
\item تاریخ منطقه را سطحی و از روی کتاب‌های درسی غربی می‌شناسند
\item ساختارهای قومی-قبیله‌ای را نمی‌فهمند یا «قرون‌وسطایی» می‌خوانند
\item دینامیک‌های اقتصادی محلی — بازار غیررسمی، شبکه‌های قاچاق،
اقتصاد معیشتی — را نادیده می‌گیرند
\item فرض می‌کنند الگوی «دولت-ملت» وستفالیایی در همه‌جا کاربرد دارد
\end{itemize}
\end{tcolorbox}

\begin{tcolorbox}[enemybox, title={عامل ۲: افق زمانی کوتاه}]
\begin{itemize}
\item چرخه‌های انتخاباتی در کشورهای غربی: ۴–۵ سال
\item زمان لازم برای ساختِ نهادهای پایدار: ۲۰–۳۰ سال (حداقل)
\item \textbf{تناقض ساختاری:} سیاست‌مداران باید «پیروزی» را زود اعلام
کنند — حتی اگر واقعیت زمینی عکس آن باشد
\item نمونه: اوباما در ۲۰۱۱ «پایان جنگ عراق» را اعلام کرد؛ سه سال
بعد داعش موصل را گرفت
\end{itemize}
\end{tcolorbox}

\begin{tcolorbox}[enemybox, title={عامل ۳: منافع متضاد}]
مداخله‌گر همزمان باید:
\begin{itemize}
\item «نجات» بدهد — وظیفه‌ی اعلام‌شده
\item منافع ملّی خود را تأمین کند — وظیفه‌ی واقعی
\item هزینه‌ها را برای مالیات‌دهندگان توجیه کند
\item افکار عمومی داخلی را راضی نگه دارد
\item متحدان منطقه‌ای خود را از دست ندهد
\end{itemize}
وقتی این اهداف با هم تعارض پیدا کنند — و \textit{همیشه} تعارض پیدا
می‌کنند — «نجات مردم» اولین قربانی است.
\end{tcolorbox}

\begin{tcolorbox}[enemybox, title={عامل ۴: وابستگی ساختاری}]
مداخله:
\begin{itemize}
\item نخبگان محلی را به مداخله‌گر وابسته می‌کند — به‌جای پاسخ‌گویی به مردم
\item اقتصاد را به «اقتصاد جنگی» تبدیل می‌کند — مترجمان، رانندگان، و
پیمان‌کاران نظامی بیشتر از معلمان درآمد دارند
\item نهادهای بومی را تضعیف و جایگزین می‌کند
\item \concept{عاملیت مردم} را می‌کُشد — مردم از «فاعل» تغییر به
«مفعولِ» آن بدل می‌شوند
\end{itemize}
\end{tcolorbox}

%% ──────────────────────────────────────────────
%%  فصل ۱۵ – بخش اضافه‌شده: «طیف مداخلات»
%%  این بخش پس از section اول (آمار تلخ) و پیش از
%%  section دوم (دست‌های کثیف) قرار می‌گیرد.
%%  فایل: chapters/ch15-paradox-of-liberation.tex (بخش ۱.۳)
%% ──────────────────────────────────────────────

\subsection{یک لحظه مکث: دقیقاً از چه سخن می‌گوییم؟}
\label{sec:ch15-taxonomy-intro}

\needspace{8\baselineskip}
پیش از آنکه در تحلیل شکست مداخلات عمیق‌تر شویم، باید به یک پرسش بنیادین
پاسخ دهیم: \textbf{وقتی می‌گوییم «مداخله‌ی خارجی»، دقیقاً به چه چیزی
اشاره می‌کنیم؟}

این پرسش از آن جهت ضروری است که در ادبیات عمومی — و حتی گاه در ادبیات
تخصصی — پدیده‌های بسیار متفاوتی زیر یک چتر قرار می‌گیرند: از استعمار
بریتانیا در هند تا حمله‌ی آمریکا به عراق، از نظام مَندیت جامعه‌ی ملل تا
بمباران ناتو در لیبی.  اینها \textbf{یک} پدیده نیستند.

\begin{tcolorbox}[defbox, title={طیف دخالت خارجی در سرنوشت ملت‌ها}]
تاریخ بشر حداقل \textbf{هشت شکل متمایز} از دخالت خارجی در سرنوشت
ملت‌ها را شناخته است:

\begin{table}[H]
\centering\small
\begin{tabularx}{\textwidth}{r|r|X|r}
\toprule
& \textbf{نوع} & \textbf{ماهیت} & \textbf{نمونه} \\
\midrule
۱ & \concept{کشورگشایی/استعمار} & تصرف و بهره‌برداری بلندمدت &
  هند بریتانیا \\
۲ & \concept{بازآرایی پساجنگی} & تقسیم غنائم + «نظم جدید» &
  مَندیت‌های عثمانی \\
۳ & \concept{استعمارزدایی} & خروج قدرت استعماری &
  تقسیم هند ۱۹۴۷ \\
۴ & \concept{مداخله‌ی ایدئولوژیک} & تغییر رژیم بر اساس ایدئولوژی &
  کودتای ۲۸ مرداد \\
۵ & \concept{مداخله‌ی بشردوستانه} & توجیه حقوق‌بشری &
  لیبی ۲۰۱۱ \\
۶ & \concept{بازیابی/وحدت‌طلبی} & الحاق به نام «هویت مشترک» &
  کریمه ۲۰۱۴ \\
۷ & \concept{مداخله‌ی ضدتروریستی} & توجیه امنیتی &
  افغانستان ۲۰۰۱ \\
۸ & \concept{مداخله‌ی اقتصادی ساختاری} & تحمیل مدل اقتصادی &
  برنامه‌های تعدیل \lat{IMF} \\
\bottomrule
\end{tabularx}
\end{table}

تحلیل مفصل این طبقه‌بندی در پیوست~\ref{ch:taxonomy-interventions}
آمده است.
\end{tcolorbox}

\begin{tcolorbox}[legalbox, title={آمار این فصل دقیقاً شامل چه می‌شود؟}]
آمار «۸۴٪ شکست» که در بخش قبل ارائه شد، عمدتاً بر پژوهش‌هایی استوار
است که \concept{تغییر رژیم تحمیلی از خارج}
(\lat{Foreign-Imposed Regime Change — FIRC}) را بررسی
کرده‌اند.\footnote{%
  Downes, Alexander B. \textit{Catastrophic Success: Why
  Foreign-Imposed Regime Change Goes Wrong}. Cornell UP, 2021;
  Peic, Goran \& Reiter, Dan. ``Foreign-Imposed Regime Change,
  State Power and Civil War Onset, 1920–2004.''
  \textit{British Journal of Political Science} 41, no.\,3 (2011).}

\textbf{شامل می‌شوند:}
\begin{itemize}
\item مداخلات نظامی ایدئولوژیک جنگ سرد (طبقه‌ی ۴)
\item مداخلات بشردوستانه (طبقه‌ی ۵)
\item مداخلات ضدتروریستی پسا ۱۱ سپتامبر (طبقه‌ی ۷)
\end{itemize}

\textbf{شامل نمی‌شوند:}
\begin{itemize}
\item استعمار کلاسیک (طبقه‌ی ۱) — پدیده‌ی متفاوت با منطق متفاوت
\item نظام مَندیت و بازآرایی پساجنگی (طبقه‌ی ۲) — «تقسیم غنائم» است،
  نه «مداخله‌ی رهایی‌بخش»
\item استعمارزدایی (طبقه‌ی ۳) — خروج است، نه ورود
\item مداخلات اقتصادی ساختاری (طبقه‌ی ۸) — معمولاً جداگانه بررسی
  می‌شوند
\end{itemize}

\textbf{چرا مهم است؟} زیرا منتقدان مداخله گاه با خلط این طبقه‌ها،
آمار را بزرگ‌نمایی می‌کنند — و مدافعان مداخله نیز با همین خلط، موارد
«موفق» را از طبقه‌ی دیگری (مثلاً بازسازی ژاپن و آلمان = طبقه‌ی ۲)
قرض می‌گیرند تا مداخلات طبقه‌ی ۵ و ۷ را توجیه کنند.
\end{tcolorbox}

این تمایز در سراسر بقیه‌ی فصل و فصل‌های آینده اهمیت دارد: وقتی از
«شکست مداخله» سخن می‌گوییم، مشخصاً به مواردی اشاره داریم که قدرتی
خارجی \textbf{با ادعای رهایی‌بخشی یا امنیت‌بخشی} وارد کشوری شده و
تلاش کرده نظام سیاسی آن را تغییر دهد. نه استعمارگری که ادعای «رهایی»
نداشت، نه استعمارزدایی که فرایندی معکوس بود.

%% ================================================================
\section{دست‌های کثیف: منتقد کیست؟}
\label{sec:ch15-dirty-hands}
%% ================================================================

\subsection{«رطب خورده کی منع رطب کند»}

\begin{histQuote}{ضرب‌المثل فارسی}
\textbf{«رطب خورده کی منع رطب کند»} —
کسی که خود از چیزی بهره برده، چگونه می‌تواند دیگران را از آن منع کند؟

این ضرب‌المثل به یکی از عمیق‌ترین معضلات نقد مداخله اشاره دارد:
\textbf{منتقدان مداخله خود اغلب از منتفعان نظم موجودند.}
\end{histQuote}

\needspace{6\baselineskip}
واقعیت آن است که:

\begin{itemize}
\item \textbf{قدرت‌های غربی} که امروز از «حقوق بشر» سخن می‌گویند، تاریخی
از استعمار، برده‌داری، و نسل‌کشی دارند — بلژیک در کنگو ۱۰ میلیون نفر را
کشت، اما امروز «درس دموکراسی» می‌دهد
\item \textbf{روسیه و چین} که از «حاکمیت ملی» دفاع می‌کنند، خود در
قلمروهایشان سرکوب‌گرند — چچن، تبت، سین‌کیانگ
\item \textbf{رژیم‌های منطقه‌ای} که از «عدم مداخله» سخن می‌گویند، خود در
کشورهای همسایه مداخله می‌کنند — ایران در سوریه و لبنان، عربستان در یمن و
بحرین، ترکیه در شمال سوریه و عراق
\item حتی \textbf{سازمان‌های حقوق بشری} که منتقد مداخله‌اند، بودجه‌شان
عمدتاً از دولت‌های غربی تأمین می‌شود
\end{itemize}


\subsection{آیا این «نقد نقد» درست است؟}

\begin{tcolorbox}[wavebox, title={تحلیل منطقی}]
\textbf{استدلال ۱:} «چون منتقد خود بی‌عیب نیست، نقدش بی‌اعتبار است»

\textbf{پاسخ:} این یک \concept{مغالطه‌ی منشأ} (\lat{Genetic Fallacy})
است. صحت یک استدلال به هویت گوینده‌اش ربطی ندارد. حتی اگر منتقد
«رطب خورده» باشد، استدلالش ممکن است درست باشد.

\tcblower

\textbf{استدلال ۲:} «قدرت‌های مداخله‌گر خود سابقه‌ی سیاه دارند»

\textbf{پاسخ:} این درست است و باید در تحلیل لحاظ شود — اما به معنای
بی‌اعتباری \textit{اصل} مداخله‌ی بشردوستانه نیست. مسئله، \textbf{اجرا}
و \textbf{انگیزه‌ها} است، نه ضرورتاً خود ایده.

\tcblower

\textbf{نتیجه:} نه منتقدان مداخله بی‌عیب‌اند، نه مدافعان آن. باید از
\textbf{هر دو} سوی بحث با دیده‌ی انتقادی نگریست.
\end{tcolorbox}


%% ================================================================
\section{شکاف درونی: چه کسی از چه کسی می‌خواهد نجات یابد؟}
\label{sec:ch15-internal-fracture}
%% ================================================================

\subsection{رهایی‌خواهی یا کسب قدرت؟}

\needspace{6\baselineskip}
یکی از مغفول‌ترین جنبه‌های بحث مداخله این است که \textbf{«درخواست نجات»
معمولاً از سوی یک زیرگروه خاص مطرح می‌شود — نه کل جامعه.}

\begin{tcolorbox}[defbox, title={پرسش‌های کلیدی}]
\begin{enumerate}
\item \textbf{چه کسی} درخواست مداخله می‌کند؟
\item این گروه \textbf{چه سهمی} از جمعیت دارد؟
\item \textbf{منافع} این گروه با منافع کل جامعه چقدر هم‌پوشانی دارد؟
\item آیا این گروه در صورت پیروزی، \textbf{حقوق سایرین} را رعایت خواهد کرد؟
\item آیا درخواست مداخله، \textbf{آخرین راه‌حل} است یا \textbf{میان‌بُرِ قدرت}؟
\item آیا اعضای این گروه حاضرند \textbf{خودشان} هزینه‌ی تغییر را بپردازند
      — یا می‌خواهند «دیگری» بجنگد؟
\end{enumerate}
\end{tcolorbox}


\subsection{نمونه‌های تاریخی}

\begin{table}[H]
\centering
\begin{tabularx}{\textwidth}{r|r|X}
\toprule
\textbf{مورد} & \textbf{گروه درخواست‌کننده} & \textbf{نتیجه برای سایر گروه‌ها} \\
\midrule
\histref{عراق}{۲۰۰۳} & بخشی از اپوزیسیون تبعیدی (عمدتاً شیعه و کرد) &
  اقلیت سنی به‌شدت به حاشیه رفت $\to$ ظهور داعش \\
\histref{لیبی}{۲۰۱۱} & شورای انتقالی (عمدتاً از بنغازی) &
  قبایل جنوب و غرب به حاشیه رفتند $\to$ جنگ داخلی \\
\histref{سوریه}{۲۰۱۱–} & گروه‌های متعدد با اهداف متضاد &
  هر گروه «نجات» خود را می‌خواست $\to$ فروپاشی \\
\histref{افغانستان}{۲۰۰۱} & ائتلاف شمال (غیرپشتون) &
  پشتون‌ها احساس طرد کردند $\to$ بازگشت طالبان \\
\histref{اوکراین}{۲۰۱۴} & جریان غرب‌گرای کی‌یف &
  شرق روسی‌زبان به حاشیه رفت $\to$ جنگ دونباس \\
\bottomrule
\end{tabularx}
\end{table}

\subsection{معضل نمایندگی}

\begin{tcolorbox}[enemybox, title={مسئله‌ی بنیادین}]
در جوامع تحت سرکوب، \textbf{هیچ مکانیزم معتبری برای سنجش خواست واقعی
مردم وجود ندارد:}
\begin{itemize}
\item انتخابات؟ $\leftarrow$ مهندسی‌شده یا ناموجود
\item نظرسنجی؟ $\leftarrow$ در فضای ترس بی‌معناست
\item اعتراضات خیابانی؟ $\leftarrow$ فقط بخشی از جامعه شرکت می‌کنند —
و اغلب بخش شهری و طبقه‌ی متوسط
\item اپوزیسیون تبعیدی؟ $\leftarrow$ سال‌ها از واقعیت جامعه دور است و
اغلب منافع کشور میزبان را بازتاب می‌دهد
\item فضای مجازی؟ $\leftarrow$ نماینده‌ی جمعیت نیست؛ اقلیت پرصدا
(\lat{loud minority}) اکثریت ساکت را نمایندگی نمی‌کند
\end{itemize}

\textbf{نتیجه:} هر ادعایی مبنی بر «نمایندگی مردم» باید با تردید جدی
مواجه شود. \thinker{ادوارد سعید} هشدار داده بود: «کسی که ادعا می‌کند
صدای بی‌صدایان است، اغلب صدای خودش را به آنها نسبت می‌دهد.»
\end{tcolorbox}


%% ================================================================
\section{رؤیای بهشت: ایده‌آلیسم رهایی‌بخش}
\label{sec:ch15-idealism}
%% ================================================================

\subsection{وقتی آرمان‌گرایی چشم را کور می‌کند}

\needspace{6\baselineskip}
بسیاری از جنبش‌های رهایی‌بخش — چه از درون، چه با کمک خارجی — بر تصویری
\concept{آرمانی} از آینده بنا شده‌اند:

\begin{itemize}
\item «بعد از سقوط دیکتاتور، دموکراسی خودبه‌خود برقرار می‌شود»
\item «مردم ذاتاً خوب‌اند؛ فقط رژیم بد است»
\item «با آزادی، اقتصاد شکوفا می‌شود»
\item «اختلافات قومی/مذهبی زیر سرکوب پنهان شده؛ با آزادی حل می‌شوند»
\item «یک انتخابات آزاد کافی‌ست تا همه‌چیز درست شود»
\end{itemize}

\begin{tcolorbox}[failbox, title={واقعیت تلخ}]
هر یک از این فرض‌ها در بیشتر موارد \textbf{نادرست} از کار درآمده است:
\begin{itemize}
\item دموکراسی نیازمند \concept{نهادها}، \concept{فرهنگ سیاسی}،
و \concept{زمان} است — آلمان غربی ۴۰ سال طول کشید تا دموکراسی‌اش
«عادی» شود
\item مردم همه خوب نیستند؛ بسیاری منتظر فرصت انتقام‌اند — رواندا نشان
داد که «مردم عادی» می‌توانند در نسل‌کشی شرکت کنند
\item اقتصاد پس از جنگ/انقلاب معمولاً \textbf{فرومی‌پاشد} — تولید
ناخالص داخلی عراق در ۲۰۰۳ نصف شد
\item اختلافات سرکوب‌شده با برداشتن درِ فشار \textbf{منفجر}
می‌شوند — یوگسلاوی پس از تیتو، عراق پس از صدام
\item اولین انتخابات پس از سقوط دیکتاتوری اغلب نه دموکرات‌ها بلکه
\textbf{سازمان‌یافته‌ترین گروه} را به قدرت می‌رساند — مثل اخوان‌المسلمین
در مصر
\end{itemize}
\end{tcolorbox}


\subsection{«چه باید کرد؟» — پرسش بی‌پاسخ؟}

اینجاست که نقد با بن‌بست مواجه می‌شود. اگر:
\begin{itemize}
\item مداخله‌ی خارجی معمولاً شکست می‌خورد
\item اپوزیسیون داخلی نماینده‌ی واقعی نیست
\item آرمان‌گرایی خطرناک است
\end{itemize}

\textbf{پس تکلیف مردم تحت ستم چیست؟}  پاسخ به این پرسش موضوع بخش‌های
بعدی این فصل و فصل‌های آینده است.


%% ================================================================
\section{معضل بنیادین: وقتی «حق» وجود ندارد}
\label{sec:ch15-no-rights}
%% ================================================================

\subsection{کسانی که از حق محروم‌اند}

\begin{histQuote}{واقعیتی تلخ}
بردگان، سیاه‌پوستان تحت آپارتاید، بومیان زیر سلطه‌ی استعمار، زنان در
نظام‌های پدرسالار شدید، اویغورها در اردوگاه‌های چین\ldots

اینها کسانی‌اند که نه فقط \textbf{حقوقشان نقض شده}، بلکه اساساً
\textbf{از داشتن حق محروم‌اند}. در نظام حقوقی حاکم، آنها «موجود» نیستند.
\thinker{هانا آرنت} این را \concept{«حق داشتن حق»}
(\lat{the right to have rights}) نامید — بنیادی‌ترین حقی که بدون آن هیچ
حق دیگری معنا ندارد.\footnote{Arendt, Hannah.
\textit{The Origins of Totalitarianism}. Harcourt, 1951, p.\,296.}
\end{histQuote}

برای این گروه‌ها، پرسش اساسی این است: \textbf{اگر نظام موجود شما را
«انسان» نمی‌شناسد، چگونه می‌توانید در چارچوب همان نظام برای حقوقتان
بجنگید؟}

\subsection{آیا «قدرت» تنها راه است؟}

\begin{tcolorbox}[failbox, title={نگاه واقع‌گرایانه}]
\begin{itemize}
\item تاریخ نشان می‌دهد که \textbf{هیچ ظالمی داوطلبانه قدرت را واگذار
نکرده}
\item بردگی در آمریکا با \textbf{جنگ داخلی}
(۶۲۰٬۰۰۰ کشته) لغو شد — نه با درخواست مؤدبانه
\item آپارتاید با \textbf{فشار بین‌المللی + مقاومت داخلی}
(دهه‌ها مبارزه) پایان یافت
\item استعمار با \textbf{جنبش‌های ضداستعماری} شکست خورد — اغلب پس از
جنگ‌های خونین (الجزایر: ۱٫۵ میلیون کشته)
\end{itemize}
\textbf{نتیجه:} بدون نوعی از «قدرت»، تغییر ممکن نیست.
\end{tcolorbox}

\begin{tcolorbox}[successbox, title={نگاه انتقادی}]
\begin{itemize}
\item \concept{قدرت کافی نیست} — نوعِ قدرت مهم است
\item قدرتی که از بیرون بیاید، \textbf{وابستگی} ایجاد می‌کند
\item قدرتی که از درون بجوشد، \textbf{پایدارتر} است
\item مهم‌تر از قدرتِ تسلیحاتی، \concept{قدرتِ ایده و روایت} است
\item حتی \thinker{گاندی} و \thinker{مارتین لوتر کینگ} از «قدرت» استفاده
کردند — اما \concept{قدرت اخلاقی}، نه نظامی
\end{itemize}
\textbf{نتیجه:} قدرت لازم است، اما \textbf{قدرت به‌تنهایی} کافی نیست.
\end{tcolorbox}


\subsection{قدرت ایده‌ها}

آنچه انقلاب‌های موفق را از ناموفق متمایز می‌کند، نه صرفاً میزان قدرت
نظامی، بلکه \concept{قدرت ایده‌ها} است:

\begin{table}[H]
\centering
\begin{tabular}{r|c|c|r}
\toprule
\textbf{انقلاب/تغییر} & \textbf{قدرت نظامی} & \textbf{قدرت ایده‌ای} &
  \textbf{نتیجه بلندمدت} \\
\midrule
\textbf{انقلاب فرانسه} & متوسط & بسیار بالا & تأثیر جهانی پایدار \\
\textbf{انقلاب آمریکا} & متوسط & بالا & دموکراسی پایدار \\
\textbf{انقلاب روسیه} & بالا & بالا & ۷۰ سال دوام \\
\textbf{انقلاب ایران} & پایین & بالا & تداوم (با چالش‌ها) \\
\textbf{تغییر رژیم عراق} & بسیار بالا & بسیار پایین & فروپاشی \\
\textbf{تغییر رژیم لیبی} & بالا & پایین & هرج‌ومرج \\
\textbf{جنبش ضدآپارتاید} & پایین & بسیار بالا & گذار موفق \\
\bottomrule
\end{tabular}
\end{table}

الگو روشن است: هرجا قدرتِ ایده‌ها بالا بود، حتی با قدرت نظامی پایین،
نتیجه‌ی بلندمدت بهتر بود. هرجا قدرت نظامی بالا بود اما ایده‌ی منسجمی
پشتش نبود، فروپاشی در انتظار بود.


%% ================================================================
\section{ایدئولوژی‌ها به‌مثابه‌ی انگیزه‌ی مداخله}
\label{sec:ch15-ideologies}
%% ================================================================

\subsection{چرا «دیگران» هزینه می‌دهند؟}

\needspace{6\baselineskip}
یک پرسش کلیدی که کمتر پرسیده می‌شود: \textbf{چرا یک قدرت خارجی باید
برای «رهایی» مردمی دیگر هزینه بدهد؟}

\begin{tcolorbox}[defbox, title={طبقه‌بندی انگیزه‌ها}]
\begin{enumerate}
\item \textbf{انگیزه‌ی ژئوپلیتیک:} تضعیف رقیب، دسترسی به منابع، کنترل
مسیرهای انرژی و بازرگانی
\item \textbf{انگیزه‌ی اقتصادی:} بازار جدید، قراردادهای بازسازی
(هالیبرتون در عراق)، نفوذ شرکتی
\item \textbf{انگیزه‌ی ایدئولوژیک:} گسترش مدل لیبرال، اسلام‌گرایی،
سوسیالیسم، پان‌ناسیونالیسم
\item \textbf{انگیزه‌ی داخلی:} انحراف توجه از بحران داخلی، اثبات
«رهبری جهانی»، رضایت افکار عمومی
\item \textbf{انگیزه‌ی نهادی:} مجتمع نظامی-صنعتی، شرکت‌های امنیتی
خصوصی، بوروکراسی‌های توسعه‌ای که بدون «بحران» بیکار می‌شوند
\item \textbf{انگیزه‌ی انسانی واقعی:} (نادر ولی موجود) فشار جامعه‌ی مدنی،
رسانه‌ها، وجدان عمومی
\end{enumerate}
\end{tcolorbox}

نکته‌ی مهم: این انگیزه‌ها \textbf{همزمان} و \textbf{درهم‌تنیده} عمل
می‌کنند. حتی صادقانه‌ترین مداخله‌ی بشردوستانه از آلودگی انگیزه‌های
دیگر مصون نیست.

\subsection{سه ایدئولوژی رهایی‌بخش}

\begin{tcolorbox}[wavebox, title={ایدئولوژی‌های پشت مداخلات «رهایی‌بخش»}]

\textbf{۱. لیبرالیسم غربی:}
\begin{itemize}
\item ادعا: گسترش دموکراسی، حقوق بشر، بازار آزاد
\item نقد: اغلب پوششی برای منافع اقتصادی و ژئوپلیتیک
\item نمونه: عراق، لیبی، «انقلاب‌های رنگی»
\item \keyword{تناقض درونی:} همان لیبرالیسم که از «رضایت حکومت‌شوندگان»
سخن می‌گوید، بدون رضایت آنها حمله می‌کند
\end{itemize}

\tcblower

\textbf{۲. اسلام‌گرایی:}
\begin{itemize}
\item ادعا: حمایت از «امت اسلامی»، مبارزه با ظلم
\item نقد: اغلب ابزار رقابت منطقه‌ای (ایران-عربستان)
\item نمونه: یمن، سوریه، بحرین
\item \keyword{تناقض درونی:} دفاع از «وحدت امت» درحالی‌که عمیق‌ترین
شکاف‌ها درون همین «امت» است
\end{itemize}

\tcblower

\textbf{۳. ناسیونالیسم قومی/نژادی:}
\begin{itemize}
\item ادعا: حمایت از «هم‌تباران» در خارج از مرزها
\item نقد: ابزار توسعه‌طلبی و بی‌ثبات‌سازی
\item نمونه: پان‌ترکیسم، پان‌عربیسم، حمایت روسیه از «روس‌زبانان»
\item \keyword{تناقض درونی:} تأکید بر «خون و خاک» درحالی‌که مرزهای
قومی هرگز تمیز نیستند
\end{itemize}
\end{tcolorbox}


\subsection{آیا ایدئولوژی ذاتاً بد است؟}

\begin{tcolorbox}[successbox, title={پاسخ کوتاه: نه}]
ایدئولوژی‌ها ابزارند — می‌توانند برای سلطه استفاده شوند، می‌توانند برای
رهایی. تفاوت در:
\begin{enumerate}
\item \textbf{چه کسی} ایدئولوژی را حمل می‌کند؟ (مردم یا قدرت‌ها)
\item \textbf{چگونه} اجرا می‌شود؟ (از پایین یا از بالا)
\item \textbf{چه فضایی} برای نقد و اصلاح وجود دارد؟ (باز یا بسته)
\item آیا \concept{عاملیت محلی} محترم شمرده می‌شود؟
\end{enumerate}
\textbf{مثال مثبت:} ایده‌های روشنگری (حقوق بشر، عقلانیت) توسط جنبش‌های
ضداستعماری در برابر \textit{خود اروپایی‌ها} به کار گرفته شد —
\thinker{توسن لوورتور} در هائیتی اعلامیه‌ی حقوق بشر فرانسه را علیه
فرانسه استفاده کرد.\footnote{James, C.\,L.\,R. \textit{The Black
Jacobins}. Vintage, 1989 [1938].}
\end{tcolorbox}


%% ================================================================
\section{نه همه‌چیز صفر است: امکان تغییر مثبت}
\label{sec:ch15-not-all-zero}
%% ================================================================

\subsection{دفاع از نومیدی؟}

\needspace{6\baselineskip}
تا اینجا ممکن است نقد ما چنین نتیجه‌ای داشته باشد: «هیچ مداخله‌ای خوب
نیست؛ هر کمکی مشکوک است؛ پس بگذارید ظلم ادامه یابد.»
\textbf{این نتیجه‌گیری نادرست و خطرناک است.}

\begin{histQuote}{سعدی شیرازی}
\textbf{بنی‌آدم اعضای یکدیگرند \quad که در آفرینش ز یک گوهرند}

\textbf{چو عضوی به درد آورد روزگار \quad دگر عضوها را نماند قرار}

این حکمت کهن ایرانی دقیقاً در برابر بی‌تفاوتی می‌ایستد. مسئله نه
\textbf{آیا} باید کمک کرد، بلکه \textbf{چگونه} و \textbf{با چه شرایطی}.
\end{histQuote}


\subsection{شرایط امکان تغییر مثبت}

\begin{tcolorbox}[successbox, title={چه زمانی کمک خارجی می‌تواند مثبت باشد؟}]
\begin{table}[H]
\centering
\begin{tabularx}{\textwidth}{r|X}
\toprule
\textbf{شرط} & \textbf{توضیح} \\
\midrule
\concept{عاملیت محلی} & جنبش اصلی از درون باشد؛ کمک خارجی فقط تقویت‌کننده \\
\concept{محدودیت زمانی} & کمک موقت و با برنامه‌ی خروج مشخص \\
\concept{عدم وابستگی} & کمک نباید وابستگی ساختاری ایجاد کند \\
\concept{شفافیت} & اهداف و منافع کمک‌کننده آشکار باشد \\
\concept{پاسخ‌گویی} & مکانیزمی برای پاسخ‌گویی کمک‌کننده به مردم محلی وجود داشته باشد \\
\concept{تناسب} & نوع کمک با نیاز واقعی متناسب باشد \\
\concept{غیرنظامی بودن} & تا حد امکان از ابزارهای غیرنظامی استفاده شود \\
\concept{فراگیری} & کمک نباید فقط یک گروه قومی/مذهبی/سیاسی را تقویت کند \\
\bottomrule
\end{tabularx}
\end{table}
\end{tcolorbox}


\subsection{نمونه‌های نسبتاً موفق}

\begin{table}[H]
\centering
\begin{tabularx}{\textwidth}{r|r|X}
\toprule
\textbf{مورد} & \textbf{نوع کمک} & \textbf{عوامل موفقیت} \\
\midrule
\textbf{آفریقای جنوبی} & تحریم + فشار دیپلماتیک &
  جنبش داخلی قوی (\lat{ANC})؛ رهبری استثنایی (ماندلا)؛ کمک غیرنظامی \\
\textbf{اروپای شرقی ۱۹۸۹} & حمایت اقتصادی و فرهنگی &
  جنبش‌های داخلی قوی؛ تغییر در مسکو؛ کمک نرم \\
\textbf{تیمور شرقی ۱۹۹۹} & مداخله‌ی سازمان ملل &
  مجوز شورای امنیت؛ رضایت اندونزی؛ اهداف محدود \\
\textbf{طرح مارشال} & بازسازی اقتصادی &
  نهادهای موجود؛ مشارکت محلی؛ هدف روشن \\
\bottomrule
\end{tabularx}
\end{table}

وجه مشترک همه‌ی این نمونه‌ها: \textbf{جنبش یا ظرفیت محلی} پیش‌تر وجود
داشت و کمک خارجی \textit{تقویت‌کننده} بود، نه \textit{جایگزین}.


%% ================================================================
\section{واکسیناسیون: چگونه از سوءاستفاده جلوگیری کنیم؟}
\label{sec:ch15-vaccination}
%% ================================================================

\subsection{مسئله‌ی «دفاع»}

یک نکته‌ی کلیدی: \textbf{طرف مدافع (رژیم موجود) نیز می‌تواند از همین
استدلال‌ها سوءاستفاده کند:}
\begin{itemize}
\item «هر مخالفتی، توطئه‌ی خارجی است»
\item «هر اعتراضی، مداخله‌ی بیگانه است»
\item «حفظ وضع موجود، دفاع از حاکمیت است»
\item «منتقدان ما عوامل استعمار نو هستند»
\end{itemize}

این دقیقاً همان «جعبه‌ی پاندورا»ی معکوس است: همان‌طور که «حق مداخله»
ابزار تجاوز می‌شود، «حق حاکمیت» هم ابزار سرکوب می‌شود.

\subsection{واکسن چیست؟}

\begin{tcolorbox}[wavebox, title={واکسیناسیون علیه هر دو نوع سوءاستفاده}]
\textbf{برای جوامع:}
\begin{enumerate}
\item \concept{تفکر انتقادی:} نپذیرفتن روایت‌های یک‌طرفه — نه از
مداخله‌گران، نه از رژیم
\item \concept{اطلاعات متنوع:} دسترسی به منابع مختلف خبری و تحلیلی
\item \concept{حافظه‌ی تاریخی:} یادآوری تجربه‌های گذشته — موفق و ناموفق
\item \concept{نهادهای مستقل:} جامعه‌ی مدنی، رسانه، دانشگاه — مستقل از
دولت \textit{و} خارج
\item \concept{گفت‌وگوی داخلی:} مکانیزم‌هایی برای حل اختلافات بدون توسل
به بیگانه
\end{enumerate}

\tcblower

\textbf{برای افراد:}
\begin{enumerate}
\item \textbf{پرسش از انگیزه‌ها:} چرا این بازیگر خارجی «کمک» می‌کند؟
\item \textbf{پرسش از هزینه‌ها:} هزینه‌ی این «کمک» چیست؟ چه کسی می‌پردازد؟
\item \textbf{پرسش از جایگزین‌ها:} آیا راه دیگری وجود ندارد؟
\item \textbf{پرسش از نتایج:} تجربه‌ی موارد مشابه چه بوده؟
\item \textbf{پرسش از خود:} آیا من هم ممکن است اشتباه کنم؟
\end{enumerate}
\end{tcolorbox}


%% ================================================================
\section{جمع‌بندی: پارادوکس حل‌نشدنی؟}
\label{sec:ch15-conclusion}
%% ================================================================

\begin{tcolorbox}[mybox, title={خلاصه‌ی فصل}]
این فصل نشان داد که:
\begin{enumerate}
\item \textbf{آمار علیه مداخله است:} $\sim$۸۴٪ شکست یا نتایج نامطلوب
\item \textbf{همه طرف‌ها دست‌های کثیف دارند:} نه منتقدان بی‌عیب‌اند، نه مدافعان
\item \textbf{رهایی‌خواهی اغلب پوشش کسب قدرت است:} گروه‌های خاص از نام «ملت»
استفاده می‌کنند
\item \textbf{آرمان‌گرایی خطرناک است:} اما بدون آرمان، تغییری ممکن نیست
\item \textbf{محرومان از حق در بن‌بست‌اند:} اما تاریخ نشان می‌دهد
راه‌هایی وجود دارد
\item \textbf{قدرت لازم است:} اما قدرت ایده‌ها مهم‌تر از قدرت نظامی است
\item \textbf{همه‌چیز صفر نیست:} امکان تغییر مثبت وجود دارد، اما با
شرایط سخت
\item \textbf{واکسیناسیون ممکن است:} تفکر انتقادی، نهادهای مستقل،
حافظه‌ی تاریخی
\end{enumerate}

\textbf{نتیجه‌ی نهایی:} پارادوکس رهایی حل‌نشدنی نیست — اما راه‌حلش ساده
هم نیست. بین بی‌تفاوتی و مداخله‌گری، راه سومی وجود دارد:
\concept{همبستگی انتقادی}. فصل‌های آینده این مفهوم را خواهند شکافت.
\end{tcolorbox}

\cleardoublepage